\Needspace{22\baselineskip}

\CNsection{1}{Why NFV?}

Traditional telecom networks rely on \textbf{proprietary hardware appliances} for every network function:

{\def\LTcaptype{none} % do not increment counter
\Needspace{14\baselineskip}
\begin{longtable}[c]{@{}
  >{\raggedright\arraybackslash\hspace{0pt}}m{(\linewidth - 2\tabcolsep) * \real{0.4091}}
  >{\raggedright\arraybackslash\hspace{0pt}}m{(\linewidth - 2\tabcolsep) * \real{0.5909}}@{}}
\toprule\noalign{}
\rowcolor{CNink}\CNth{Problem} & \CNth{Description} \\
\CNheadrulesep
\endhead
\bottomrule\noalign{}
\endlastfoot
\textbf{High CAPEX} & Each function requires dedicated, expensive hardware (firewall appliance, load balancer, DPI box, etc.) \\
\textbf{High OPEX} & Power, cooling, rack space, vendor maintenance contracts for each appliance \\
\textbf{Slow Scaling} & Need more capacity? Order hardware, wait weeks for delivery, rack-and-stack, configure \\
\textbf{Underutilization} & Hardware sized for peak load; average utilization often 20-30\% \\
\textbf{Vendor Dependency} & Locked into vendor hardware lifecycle, upgrade timelines, and pricing \\
\textbf{Rigid Deployment} & Cannot easily move functions, resize, or replicate on demand \\
\end{longtable}
}

\begin{CNquote}

\textbf{The Core Insight}: If network functions (firewall, router, load balancer) are just software, why not run them on standard servers as VMs or containers?

\end{CNquote}

\textbf{NFV was initiated by a group of telecom operators (AT\&T, BT, Deutsche Telekom, etc.) who published the ETSI NFV white paper in 2012.}

\Needspace{14\baselineskip}

\CNkeepfig{m16-traditional}{8}

\CNsection{2}{Traditional vs NFV Comparison}

\CNchained\CNsubsection{}{Traditional Approach}

\CNfigure{m16-traditional}{Traditional approach: one proprietary appliance per function}

\CNkeepfig{m16-nfv-approach}{3}

\CNsubsection{}{NFV Approach}

\CNfigure{m16-nfv-approach}{NFV approach: virtual functions on shared commodity hardware}

{\def\LTcaptype{none} % do not increment counter
\Needspace{14\baselineskip}
\begin{longtable}[c]{@{}
  >{\raggedright\arraybackslash\hspace{0pt}}m{(\linewidth - 4\tabcolsep) * \real{0.3077}}
  >{\raggedright\arraybackslash\hspace{0pt}}m{(\linewidth - 4\tabcolsep) * \real{0.5000}}
  >{\raggedright\arraybackslash\hspace{0pt}}m{(\linewidth - 4\tabcolsep) * \real{0.1923}}@{}}
\toprule\noalign{}
\rowcolor{CNink}\CNth{Aspect} & \CNth{Traditional} & \CNth{NFV} \\
\CNheadrulesep
\endhead
\bottomrule\noalign{}
\endlastfoot
Hardware & Proprietary, purpose-built & COTS (Commercial Off-The-Shelf) servers \\
Software & Embedded in appliance & Decoupled, runs as VM/\allowbreak{}container \\
Scaling & Buy more hardware & Spin up more instances (minutes) \\
Cost & High CAPEX + OPEX & Lower (commodity HW + SW licenses) \\
Flexibility & Fixed function & Any function on any server \\
Lifecycle & 5-7 year hardware refresh & Continuous software updates \\
\end{longtable}
}

\CNkeepfig{m16-etsi-nfv}{5}

\CNsection{3}{ETSI NFV Architecture}

\CNfigure{m16-etsi-nfv}{ETSI NFV reference architecture}

\Needspace{13\baselineskip}

\CNsubsection{}{Three Main Components:}

{\def\LTcaptype{none} % do not increment counter
\Needspace{9\baselineskip}
\begin{longtable}[c]{@{}
  >{\raggedright\arraybackslash\hspace{0pt}}m{(\linewidth - 2\tabcolsep) * \real{0.6471}}
  >{\raggedright\arraybackslash\hspace{0pt}}m{(\linewidth - 2\tabcolsep) * \real{0.3529}}@{}}
\toprule\noalign{}
\rowcolor{CNink}\CNth{Component} & \CNth{Role} \\
\CNheadrulesep
\endhead
\bottomrule\noalign{}
\endlastfoot
\textbf{NFVI} & The infrastructure: hardware + virtualization layer that hosts VNFs \\
\textbf{VNF} & Software implementation of a network function running on NFVI \\
\textbf{MANO} & Management and Orchestration: lifecycle, resources, service chaining \\
\end{longtable}
}

\Needspace{20\baselineskip}

\CNsection{4}{NFV Infrastructure (NFVI)}

NFVI provides the environment where VNFs execute.

\CNsubsection{4.1}{Hardware Resources}

{\def\LTcaptype{none} % do not increment counter
\Needspace{9\baselineskip}
\begin{longtable}[c]{@{}cc@{}}
\toprule\noalign{}
\rowcolor{CNink}\CNth{Resource} & \CNth{Description} \\
\CNheadrulesep
\endhead
\bottomrule\noalign{}
\endlastfoot
\textbf{Compute} & x86/\allowbreak{}ARM CPUs, GPU accelerators, NUMA topology \\
\textbf{Storage} & Local SSD/\allowbreak{}NVMe, SAN, distributed storage (Ceph) \\
\textbf{Network} & NICs (10/\allowbreak{}25/\allowbreak{}100G), SmartNICs, switches, SR-IOV capable \\
\end{longtable}
}

\Needspace{13\baselineskip}

\CNsubsection{4.2}{Virtualization Layer}

{\def\LTcaptype{none} % do not increment counter
\Needspace{9\baselineskip}
\begin{longtable}[c]{@{}
  >{\raggedright\arraybackslash\hspace{0pt}}m{(\linewidth - 2\tabcolsep) * \real{0.4800}}
  >{\raggedright\arraybackslash\hspace{0pt}}m{(\linewidth - 2\tabcolsep) * \real{0.5200}}@{}}
\toprule\noalign{}
\rowcolor{CNink}\CNth{Technology} & \CNth{Description} \\
\CNheadrulesep
\endhead
\bottomrule\noalign{}
\endlastfoot
\textbf{Hypervisor} & KVM, VMware ESXi, Xen — creates VMs \\
\textbf{Container Runtime} & Docker, containerd, CRI-O — creates containers \\
\textbf{Kubernetes} & Container orchestration platform (for CNFs) \\
\end{longtable}
}

\CNsubsection{4.3}{COTS Servers}

\begin{itemize}
\tightlist
\item
  \textbf{Commercial Off-The-Shelf} servers (Dell, HP, Supermicro)
\item
  Standard x86 architecture — no proprietary hardware
\item
  High-volume, competitive pricing
\item
  Performance enhanced with: SR-IOV, DPDK, huge pages, CPU pinning, NUMA awareness
\end{itemize}

\Needspace{17\baselineskip}

\CNsubsection{4.4}{Acceleration Technologies}

{\def\LTcaptype{none} % do not increment counter
\Needspace{13\baselineskip}
\begin{longtable}[c]{@{}cc@{}}
\toprule\noalign{}
\rowcolor{CNink}\CNth{Technology} & \CNth{Purpose} \\
\CNheadrulesep
\endhead
\bottomrule\noalign{}
\endlastfoot
\textbf{SR-IOV} & Bypass hypervisor for direct NIC access (low latency) \\
\textbf{DPDK} & User-space packet processing (bypass kernel) \\
\textbf{SmartNIC/\allowbreak{}DPU} & Offload networking to NIC hardware \\
\textbf{FPGA} & Programmable hardware acceleration \\
\textbf{Huge Pages} & Reduce TLB misses for high-throughput workloads \\
\end{longtable}
}

\CNkeepfig{m16-vnf-vs-cnf}{5}

\CNsection{5}{VNF vs CNF}

\CNfigure{m16-vnf-vs-cnf}{VNF (virtual machines) versus CNF (containers)}

\Needspace{25\baselineskip}

\CNsubsection{}{Comparison}

{\def\LTcaptype{none} % do not increment counter
\Needspace{21\baselineskip}
\begin{longtable}[c]{@{}
  >{\raggedright\arraybackslash\hspace{0pt}}m{(\linewidth - 4\tabcolsep) * \real{0.1739}}
  >{\raggedright\arraybackslash\hspace{0pt}}m{(\linewidth - 4\tabcolsep) * \real{0.3478}}
  >{\raggedright\arraybackslash\hspace{0pt}}m{(\linewidth - 4\tabcolsep) * \real{0.4783}}@{}}
\toprule\noalign{}
\rowcolor{CNink}\CNth{Aspect} & \CNth{VNF (VM-based)} & \CNth{CNF (Container-based)} \\
\CNheadrulesep
\endhead
\bottomrule\noalign{}
\endlastfoot
\textbf{Runtime} & Hypervisor (KVM, VMware) & Container runtime (K8s) \\
\textbf{Size} & GBs (full OS image) & MBs (shared kernel) \\
\textbf{Boot Time} & Minutes & Seconds \\
\textbf{Architecture} & Monolithic & Microservices \\
\textbf{State} & Stateful (local state in VM) & Stateless (external state store) \\
\textbf{Scaling} & Vertical (bigger VM) or slow horizontal & Fast horizontal (replicas) \\
\textbf{Updates} & Disruptive (VM restart) & Rolling updates, canary \\
\textbf{Lifecycle} & VNFM manages & Kubernetes manages (Helm, Operators) \\
\textbf{Density} & Low (few VMs per server) & High (hundreds of containers) \\
\textbf{CI/\allowbreak{}CD} & Difficult & Native (GitOps, pipelines) \\
\end{longtable}
}

\CNsubsection{}{5G Core = CNFs}

The 3GPP 5G Core (5GC) is designed as a \textbf{Service-Based Architecture (SBA)} with cloud-native principles:

\begin{itemize}
\tightlist
\item
  Each Network Function (AMF, SMF, UPF, NRF, etc.) is a \textbf{CNF}
\item
  Deployed on \textbf{Kubernetes}
\item
  Communicates via \textbf{HTTP/\allowbreak{}2 + REST APIs}
\item
  Stateless design with external data stores (e.g., UDSF)
\item
  Horizontal scaling, rolling upgrades, self-healing
\end{itemize}

\Needspace{14\baselineskip}

\Needspace{22\baselineskip}

\CNsection{6}{MANO (Management and Orchestration)}

\CNchained\CNsubsection{6.1}{NFVO (NFV Orchestrator)}

{\def\LTcaptype{none} % do not increment counter
\Needspace{13\baselineskip}
\begin{longtable}[c]{@{}
  >{\raggedright\arraybackslash\hspace{0pt}}m{(\linewidth - 2\tabcolsep) * \real{0.4348}}
  >{\raggedright\arraybackslash\hspace{0pt}}m{(\linewidth - 2\tabcolsep) * \real{0.5652}}@{}}
\toprule\noalign{}
\rowcolor{CNink}\CNth{Function} & \CNth{Description} \\
\CNheadrulesep
\endhead
\bottomrule\noalign{}
\endlastfoot
\textbf{Service Orchestration} & Composes end-to-end network services from multiple VNFs \\
\textbf{Lifecycle Management} & Onboard, instantiate, scale, update, terminate services \\
\textbf{Resource Coordination} & Requests resources from VIM across multiple data centers \\
\textbf{Catalog Management} & Stores NSD (Network Service Descriptors) and VNF packages \\
\textbf{Policy Management} & Placement, affinity, anti-affinity rules \\
\end{longtable}
}

Examples: ONAP, OSM (Open Source MANO), Cloudify, Nokia CloudBand

\Needspace{18\baselineskip}

\CNsubsection{6.2}{VNFM (VNF Manager)}

{\def\LTcaptype{none} % do not increment counter
\Needspace{14\baselineskip}
\begin{longtable}[c]{@{}
  >{\raggedright\arraybackslash\hspace{0pt}}m{(\linewidth - 2\tabcolsep) * \real{0.4348}}
  >{\raggedright\arraybackslash\hspace{0pt}}m{(\linewidth - 2\tabcolsep) * \real{0.5652}}@{}}
\toprule\noalign{}
\rowcolor{CNink}\CNth{Function} & \CNth{Description} \\
\CNheadrulesep
\endhead
\bottomrule\noalign{}
\endlastfoot
\textbf{Instantiate} & Deploy VNF from package (image + descriptor) \\
\textbf{Scale} & Scale out (add instances) or scale in (remove) \\
\textbf{Heal} & Detect failure and restart/\allowbreak{}replace VNF instances \\
\textbf{Terminate} & Graceful shutdown and resource release \\
\textbf{Update/\allowbreak{}Upgrade} & Apply new software version \\
\textbf{Configure} & Day-0 (initial), Day-1 (service), Day-2 (operational) config \\
\end{longtable}
}

\Needspace{15\baselineskip}

\CNsubsection{6.3}{VIM (Virtualized Infrastructure Manager)}

{\def\LTcaptype{none} % do not increment counter
\Needspace{11\baselineskip}
\begin{longtable}[c]{@{}
  >{\raggedright\arraybackslash\hspace{0pt}}m{(\linewidth - 2\tabcolsep) * \real{0.4348}}
  >{\raggedright\arraybackslash\hspace{0pt}}m{(\linewidth - 2\tabcolsep) * \real{0.5652}}@{}}
\toprule\noalign{}
\rowcolor{CNink}\CNth{Function} & \CNth{Description} \\
\CNheadrulesep
\endhead
\bottomrule\noalign{}
\endlastfoot
\textbf{Compute Management} & Create/\allowbreak{}delete VMs or containers, CPU/\allowbreak{}memory allocation \\
\textbf{Network Management} & Virtual networks, subnets, ports, security groups \\
\textbf{Storage Management} & Volumes, snapshots, object storage \\
\textbf{Resource Monitoring} & Utilization, capacity, alarms \\
\end{longtable}
}

Examples:

\begin{itemize}
\tightlist
\item
  \textbf{OpenStack} — most common VIM for VNFs (Nova, Neutron, Cinder)
\item
  \textbf{Kubernetes} — VIM for CNFs (native container orchestration)
\item
  \textbf{VMware vCenter} — enterprise virtualization
\end{itemize}

\Needspace{14\baselineskip}

\Needspace{20\baselineskip}

\CNsection{7}{NFV in Telecom (Real Examples)}

\CNchained\CNsubsection{7.1}{vEPC (Virtualized Evolved Packet Core)}

{\def\LTcaptype{none} % do not increment counter
\Needspace{11\baselineskip}
\begin{longtable}[c]{@{}ccc@{}}
\toprule\noalign{}
\rowcolor{CNink}\CNth{Component} & \CNth{Traditional} & \CNth{Virtualized} \\
\CNheadrulesep
\endhead
\bottomrule\noalign{}
\endlastfoot
MME & Dedicated chassis & VM running MME software \\
S-GW & Proprietary appliance & VM/\allowbreak{}container with S-GW logic \\
P-GW & Proprietary appliance & VM/\allowbreak{}container with P-GW logic \\
HSS & Dedicated server & VM with database \\
\end{longtable}
}

Benefits: scale MME independently during mass events, deploy regional instances

\CNsubsection{7.2}{vIMS (Virtualized IP Multimedia Subsystem)}

\begin{itemize}
\tightlist
\item
  P-CSCF, S-CSCF, I-CSCF as VNFs
\item
  Scales voice/\allowbreak{}video capacity on demand
\item
  Deploy in new markets without shipping hardware
\end{itemize}

\CNkeepfig{m16-k8s-5gc}{3}

\CNsubsection{7.3}{5G Core as CNFs on Kubernetes}

\CNfigure{m16-k8s-5gc}{5G core functions as CNFs on Kubernetes}

\begin{itemize}
\tightlist
\item
  Each 5G NF is a Kubernetes deployment with multiple replicas
\item
  Helm charts or Operators manage lifecycle
\item
  GitOps pipelines for CI/\allowbreak{}CD
\item
  Horizontal Pod Autoscaler for dynamic scaling
\end{itemize}

\CNkeepfig{m16-evolution}{5}

\CNsection{8}{Evolution: Physical \ensuremath{\to} Virtual \ensuremath{\to} Cloud-Native}

\CNfigure{m16-evolution}{From physical appliances to cloud-native functions}

{\def\LTcaptype{none} % do not increment counter
\Needspace{9\baselineskip}
\begin{longtable}[c]{@{}
  >{\raggedright\arraybackslash\hspace{0pt}}m{(\linewidth - 8\tabcolsep) * \real{0.2069}}
  >{\raggedright\arraybackslash\hspace{0pt}}m{(\linewidth - 8\tabcolsep) * \real{0.2241}}
  >{\raggedright\arraybackslash\hspace{0pt}}m{(\linewidth - 8\tabcolsep) * \real{0.2586}}
  >{\raggedright\arraybackslash\hspace{0pt}}m{(\linewidth - 8\tabcolsep) * \real{0.1552}}
  >{\raggedright\arraybackslash\hspace{0pt}}m{(\linewidth - 8\tabcolsep) * \real{0.1552}}@{}}
\toprule\noalign{}
\rowcolor{CNink}\CNth{Generation} & \CNth{Form Factor} & \CNth{Orchestration} & \CNth{Scaling} & \CNth{Updates} \\
\CNheadrulesep
\endhead
\bottomrule\noalign{}
\endlastfoot
\textbf{Physical} & Proprietary appliance & Manual /\allowbreak{} NMS & Buy hardware & Firmware upgrade (downtime) \\
\textbf{Virtual (VNF)} & VM on hypervisor & MANO (NFVO/\allowbreak{}VNFM) & New VMs (minutes) & VM image replacement \\
\textbf{Cloud-Native (CNF)} & Container on K8s & Kubernetes + GitOps & Pod replicas (seconds) & Rolling /\allowbreak{} Canary /\allowbreak{} Blue-Green \\
\end{longtable}
}

\CNsubsection{}{Why the Industry Moved to CNF:}

\begin{enumerate}
\def\labelenumi{\arabic{enumi}.}
\tightlist
\item
  \textbf{Faster scaling}: Seconds vs. minutes
\item
  \textbf{Higher density}: More functions per server
\item
  \textbf{DevOps native}: CI/\allowbreak{}CD, GitOps, automated testing
\item
  \textbf{Resilience}: Self-healing, stateless design
\item
  \textbf{Portability}: Run anywhere Kubernetes runs (on-prem, public cloud, edge)
\end{enumerate}

\Needspace{26\baselineskip}

\CNsection{9}{Benefits of NFV}

{\def\LTcaptype{none} % do not increment counter
\Needspace{20\baselineskip}
\begin{longtable}[c]{@{}
  >{\raggedright\arraybackslash\hspace{0pt}}m{(\linewidth - 2\tabcolsep) * \real{0.4091}}
  >{\raggedright\arraybackslash\hspace{0pt}}m{(\linewidth - 2\tabcolsep) * \real{0.5909}}@{}}
\toprule\noalign{}
\rowcolor{CNink}\CNth{Benefit} & \CNth{Description} \\
\CNheadrulesep
\endhead
\bottomrule\noalign{}
\endlastfoot
\textbf{Flexibility} & Deploy any function on any server; relocate workloads \\
\textbf{Elastic Scaling} & Scale out/\allowbreak{}in based on demand (auto-scaling) \\
\textbf{Multi-tenancy} & Multiple tenants share infrastructure securely \\
\textbf{CI/\allowbreak{}CD} & Continuous integration and deployment of network functions \\
\textbf{Cost Reduction} & 40-60\% CAPEX reduction, 30-50\% OPEX reduction (industry estimates) \\
\textbf{Speed to Market} & New services in days/\allowbreak{}weeks vs. months/\allowbreak{}years \\
\textbf{Vendor Diversity} & Mix VNFs from different vendors on same infrastructure \\
\textbf{Geographic Flexibility} & Deploy at edge, regional DC, or central cloud \\
\textbf{Resource Efficiency} & Better utilization through sharing and dynamic allocation \\
\end{longtable}
}

\Needspace{14\baselineskip}

\CNsection{10}{Challenges of NFV}

\CNchained\CNsubsection{10.1}{Performance (Packet Processing)}

\begin{itemize}
\tightlist
\item
  VMs/\allowbreak{}containers add overhead vs. bare-metal ASIC processing
\item
  \textbf{Mitigation}: DPDK, SR-IOV, SmartNICs, CPU pinning, huge pages
\item
  UPF and other data-plane VNFs are most sensitive
\end{itemize}

\CNsubsection{10.2}{Orchestration Complexity}

\begin{itemize}
\tightlist
\item
  MANO stack is complex to deploy and operate
\item
  Multi-vendor VNF interoperability issues
\item
  Descriptor standardization still evolving (TOSCA, SOL001/\allowbreak{}006)
\item
  Day-2 operations (monitoring, healing, scaling policies) are hard
\end{itemize}

\CNsubsection{10.3}{State Management}

\begin{itemize}
\tightlist
\item
  VNFs often designed as stateful monoliths (ported from appliances)
\item
  Stateful VMs are hard to scale horizontally or migrate
\item
  CNF approach: externalize state (Redis, etcd, databases)
\item
  Challenge: latency of external state access for high-throughput functions
\end{itemize}

\CNsubsection{10.4}{Reliability and Availability}

\begin{itemize}
\tightlist
\item
  Software failures more frequent than hardware failures
\item
  Need robust health monitoring, auto-healing, redundancy
\item
  Carrier-grade availability (99.999\%) requires careful design
\end{itemize}

\CNsubsection{10.5}{Security}

\begin{itemize}
\tightlist
\item
  Larger attack surface (hypervisor, orchestration APIs, shared infrastructure)
\item
  Multi-tenant isolation must be enforced at multiple layers
\item
  Container escape vulnerabilities
\end{itemize}

\CNsubsection{10.6}{Skills Gap}

\begin{itemize}
\tightlist
\item
  Operators must learn cloud technologies (OpenStack, K8s, CI/\allowbreak{}CD)
\item
  Telecom + IT convergence requires new organizational models
\end{itemize}

\CNboxsection{Key Takeaways}

\begin{CNbox}[unbreakable]{sum}{Key Takeaways}

\begin{enumerate}
\def\labelenumi{\arabic{enumi}.}
\tightlist
\item
  NFV \textbf{decouples network functions from proprietary hardware} \ensuremath{\to} software on COTS
\item
  ETSI NFV architecture: \textbf{NFVI + VNF + MANO}
\item
  \textbf{VNFs} (VM-based) are being replaced by \textbf{CNFs} (container-based, cloud-native)
\item
  \textbf{5G Core is designed as CNFs} on Kubernetes from the start
\item
  \textbf{MANO} orchestrates the full lifecycle: instantiate, scale, heal, terminate
\item
  Evolution: Physical \ensuremath{\to} Virtual (VNF) \ensuremath{\to} \textbf{Cloud-Native (CNF)}
\item
  Key challenges: \textbf{performance, orchestration complexity, and state management}
\item
  NFV + SDN together enable \textbf{fully programmable, automated telecom networks}
\end{enumerate}

\end{CNbox}
